\documentclass[10pt]{article}

\usepackage[margin=0.75in]{geometry}
\usepackage{amsmath,amsthm,amssymb,bm}
\usepackage{xcolor}
\usepackage{cancel}
\usepackage{graphicx}
\usepackage{changepage}
\usepackage{circuitikz}
\usepackage{pgfplots}
\usepackage{minted}
\usepackage{physics}
\usepackage{tikz}
\usepackage{tikz-3dplot}
\usepackage{hyperref}
\usepackage{cleveref}
\usepackage{siunitx}
\usepackage{fontspec}
\usepackage{relsize}
\usepackage{subfig}
\usepackage{todonotes}
\usepackage{contour}
\usepackage{pgffor}
\usepackage{multicol, multirow, booktabs}
\usepackage[breakable]{tcolorbox}
\usepackage[inline]{enumitem}

\theoremstyle{definition}
\newtheorem{problem}{Problem}
\newtheorem{soln}{Solution}

\pgfplotsset{compat=newest}
\usetikzlibrary{lindenmayersystems}
\usetikzlibrary{arrows}
\usetikzlibrary{calc}
\usetikzlibrary{positioning, fit}
\usetikzlibrary{3d, perspective}
\usetikzlibrary{math}
\usetikzlibrary{fadings}
\usetikzlibrary{angles,quotes} 
\usetikzlibrary{decorations.markings}

\definecolor{incolor}{HTML}{303F9F}
\definecolor{outcolor}{HTML}{D84315}
\definecolor{cellborder}{HTML}{CFCFCF}
\definecolor{cellbackground}{HTML}{F7F7F7}
\newcommand{\ui}{\hat{i}}
\newcommand{\uj}{\hat{j}}
\newcommand{\uk}{\hat{k}}
\newcommand{\ux}{\hat{\mathbf{x}}}
\newcommand{\uy}{\hat{\mathbf{y}}}
\newcommand{\uz}{\hat{\mathbf{z}}}
\newcommand{\uv}[1]{\hat{\bv{{#1}}}}
\newcommand{\pr}[1]{{#1^\prime}}
\newcommand{\pd}[2][]{{\frac{\partial^{#1}}{\partial {{#2}^{#1}}}}}
\newcommand{\dif}[2][]{{\frac{d^{#1}}{d{{#2}^{#1}}}}}
\newcommand{\grd}{\nabla}

\pgfdeclarelayer{background}  
\pgfsetlayers{background,main}
\AtBeginDocument{\RenewCommandCopy\qty\SI}

\makeatletter
\newcommand{\boxspacing}{\kern\kvtcb@left@rule\kern\kvtcb@boxsep}
\makeatother
\newcommand{\prompt}[4]{
    \ttfamily\llap{{\color{#2}[#3]:\hspace{3pt}#4}}\vspace{-\baselineskip}
}

\newcommand{\thevenin}[2]{
  \begin{center}
    \begin{circuitikz} \draw
      (0,0) -- (2,0) to[battery1, l_=$V_{Th}\eq#1$] (2,2) 
      to[resistor, l_=$R_{Th}\eq#2$] (0,2)
      ;
      \draw [o-] (-.07,2.079);
      \draw [o-] (-.07,0.079);
    \end{circuitikz}
  \end{center}
}

\newcommand{\norton}[2]{
  \begin{center}
    \begin{circuitikz} \draw
      (0,0) -- (3,0) to[american current source, l_=$I_{N}\eq#1$] (3,2) -- (0,2) (2,0)
      to[resistor, l=$R_{N}\eq#2$] (2,2)
      ;
      \draw [o-] (-.07,2.079);
      \draw [o-] (-.07,0.079);
    \end{circuitikz}
  \end{center}
}

\DeclareMathOperator{\Div}{div}
\DeclareMathOperator{\Curl}{curl}
\DeclareMathOperator{\sgn}{sgn}

\newcommand{\highlight}[1]{\colorbox{yellow}{$\displaystyle #1$}}
\newcommand{\ti}[1]{\widetilde{#1}}
\renewcommand{\div}[1]{\bv{\nabla}\cdot #1}
\renewcommand{\curl}[1]{\bv{\nabla}\times #1}

\newfontface{\Kaufmann}{Kaufmann}
\DeclareTextFontCommand{\kf}{\Kaufmann}
\newcommand{\scriptr}{\fontsize{12pt}{12pt}\kf{r}\fontsize{10pt}{10pt}}

\newfontface{\KaufmannB}{Kaufmann Bold}
\DeclareTextFontCommand{\kfb}{\KaufmannB}
\newcommand{\bscriptr}{\fontsize{12pt}{12pt}\kfb{r}\fontsize{10pt}{10pt}}
\newcommand{\justif}[2]{&{#1}&\text{#2}}
\newcommand{\bv}[1]{\bm{\mathrm{#1}}}

\title{Physics 4140H: Quiz II}
\author{Jeremy Favro (0805980) \\ Trent University, Peterborough, ON, Canada}
\date{\today}

\begin{document}
\maketitle

% PROBLEM 1
\begin{problem}
Assume that the cost of flying an aircraft at height $z$ is $e^{-kz}$ per unit distance of flight-path, where $k$ is a
positive constant. Consider that the aircraft flies in the $(x, z)$-plane from the point $(-a, 0)$ to the point $(a,
  0)$ where $z = 0$ corresponds to ground level, and where the $z$-axis points vertically upwards. Find the
extrema for the problem. The cost integral to be minimized is
$$\int_{-a}^{a}e^{-kz}\sqrt{1+(z')^2}dz.$$
Show that
$$z(x)=\frac{1}{k}\ln(\frac{\cos(kx)}{\cos(ka)}).$$
Use the boundary conditions given above to arrive at the answer.
\end{problem}
\begin{soln}
  Applying the Euler-Lagrange equation for $f=e^{-kz}\sqrt{1+(z')^2}$,
  \begin{align*}
    \frac{\partial f}{\partial z}                & =-ke^{-kz}\sqrt{1+(z')^2}                                                                                                         \\
    \frac{\partial f}{\partial z'}               & =e^{-kz}\frac{z'}{\sqrt{1+(z')^2}}                                                                                                \\
    \frac{d}{dx}(\frac{\partial f}{\partial z'}) & =-ke^{-kz}\frac{(z')^2}{\sqrt{1+(z')^2}}+e^{-kz}\left[\frac{z''\sqrt{1+(z')^2}-\frac{(z')^2z''}{\sqrt{1+(z')^2}}}{1+(z')^2}\right] \\
                                                 & =e^{-kz}\left[-\frac{k(z')^2}{\sqrt{1+(z')^2}}+\frac{z''}{\sqrt{1+(z')^2}}-\frac{(z')^2z''}{\left[1+(z')^2\right]^{3/2}}\right]
  \end{align*}
  which when plugged into the EL equation gives
  $$-ke^{-kz}\sqrt{1+(z')^2}=e^{-kz}\left[-\frac{k(z')^2}{\sqrt{1+(z')^2}}+\frac{z''}{\sqrt{1+(z')^2}}-\frac{(z')^2z''}{\left[1+(z')^2\right]^{3/2}}\right]$$
  from which we can cancel the $e^-kz$ terms as they are never zero, then multiply by $\sqrt{1+(z')^2}$ to obtain
  $$-k-\cancel{k(z')^2}=-\cancel{k(z')^2}+z''-\frac{(z')^2z''}{1+(z')^2}\implies z''-\frac{(z')^2z''}{1+(z')^2}+k=0.$$
  We then multiply through by the denominator of the remaining fraction to obtain
  $$z''(1+(z')^2)-(z')^2z''+k(1+(z')^2)=0 \implies z''+\cancel{z''(z')^2}-\cancel{(z')^2z''}+k+k(z')^2=0$$
  which finally yields the expression
  $$z''+k(1+(z')^2)=0.$$
  Now solving this took me a while, I went way off the right path at first and got stuck at a really horrible integral. The correct way to do this, I think (at least it gets me to the right answer)
  is to rewrite the primes in Leibnitz notation,
$$\frac{d}{dx}\frac{dz}{dx}+k\left(1+\left(\frac{dz}{dx}\right)^2\right)=0$$
rearrange
$$\frac{d}{dx}\frac{dz}{dx}\frac{1}{1+\left(\frac{dz}{dx}\right)^2}=-k$$
multiply by $dx$ and integrate
$$\int\frac{1}{1+\left(\frac{dz}{dx}\right)^2}d\frac{dz}{dx}=-\int kdx$$
and obtain
$$\int\frac{1}{1+\left(\frac{dz}{dx}\right)^2}d\frac{dz}{dx}=-kx+C.$$
The integral of $1/(1+x^2)$ is doable by hand but we can get from a table that it comes out to $\arctan(x)$ so 
$$\arctan(\frac{dz}{dx})=-kx+C\implies \tan(C-kx)dx=dz\implies z(x)=-\frac{1}{k}\ln(\cos(C-kx))$$
where I've used the fact that the integral of $\tan(x)$ is $\ln(\sec(x))=-\ln(\cos(x))$. Now we evaluate this between
$(-a,0)$ and $(x,0)$ (where I've flipped the order because of the negative so we can get rid of it) to obtain, using $\ln(a)-\ln(b)=\ln(a/b)$,
$$z(x)=\frac{1}{k}\ln(\frac{\cos(C-kx)}{\cos(C+ka)}).$$
The boundaries require that $z(\pm a)=\frac{1}{k}\ln(\frac{\cos(C\pm ka)}{\cos(C+ka)})=0\implies C=0$ as $\ln(1)=0$.
\end{soln}
\end{document}